Indeed, let @@M0@@ be the generic point of @@M1@@ and @@M2@@ the algebraic closure of @@M3@@. It is clear that @@M4@@ is the limit of the filtered inverse system @@M5@@ of integral affine schemes endowed with a finite faithfully flat morphism @@M6@@, where @@M7@@ is a nonempty open subset of @@M8@@. The normalization of @@M9@@ therefore comes from an @@M10@@-scheme @@M11@@ for some @@M12@@, and the finite surjective morphism @@M13@@ from a finite surjective @@M14@@-morphism @@M15@@ (EGA IV 8.8.2, 8.10.5). Finally, since @@M16@@ is normal and integral, there exists an open subset @@M17@@ of @@M18@@ such that @@M19@@ is normal and integral over @@M20@@ (EGA IV 9.9.5).

\sgaStableTarget{sga6.e13.rem.3-10}\textbf{Remark 3.10.} Retaining the notation and conditions of 3.9, suppose that there exists a resolving morphism @@M21@@ (EGA IV 7.9.1), (for example if @@M22@@, by Abhyankar). Then a similar argument shows that there exist a finite faithfully flat morphism @@M23@@, a scheme @@M24@@ smooth over @@M25@@ and a proper surjective birational morphism @@M26@@.

In the proof of the first part of 3.8, we reduced to the case where @@M27@@ is integral and @@M28@@ is normal and integral over @@M29@@, so @@M30@@ is integral. If @@M31@@, then @@M32@@ and the assertion is trivial. Otherwise, after restricting @@M33@@, one may suppose that for every @@M34@@ we have @@M35@@, so @@M36@@ is @@M37@@-regular and the dimension @@M38@@ of @@M39@@ is @@M40@@.

Now let @@M41@@ be a bounded family of classes of invertible sheaves on the fibers of @@M42@@. Let @@M43@@ be an invertible sheaf on a geometric fiber @@M44@@ of @@M45@@, with Hilbert polynomial @@M46@@; then the Hilbert polynomial of @@M47@@ is @@M48@@. As the class of @@M49@@ ranges over @@M50@@, that is, the
